% Source: Wald GR, physical PDF pp. 224--225
%         (printed pp. 211--212).
% Coverage: Chapter 9 opening material before Section 9.1.
% Figures/tables/footnotes: none.
% Status: source-reviewed and compile-clean.
% Uncertainties: none.

% Wald numbers results within each subsection (e.g., Theorem 9.2.1) and restarts
% footnotes at the opening of the chapter.
\begingroup
\edef\chapterninepreviousfootnote{\number\value{footnote}}
\setcounter{footnote}{0}
\renewcommand{\thetheorem}{\thesubsection.\arabic{theorem}}

In chapter~5 we studied the dynamics of a homogeneous, isotropic universe. We
found there that if the universe presently is expanding (as it is observed to be), then
a finite time ago it must have begun in a singular state, with infinite density and
infinite spacetime curvature. In section~6.4 we analyzed the extended Schwarzschild
geometry. There we found that at ``\(r=0\)'' an infinite spacetime curvature singularity
is present. Thus, a singularity also is predicted in the complete gravitational collapse
of a spherical body.

However, the above predictions of singularities in cosmology and gravitational
collapse are based solely on the analysis of solutions with a very high degree of
symmetry. It certainly is possible that they could give a completely misleading
picture of singularity formation. For example, in the Newtonian theory of gravity,
if a spherical, nonrotating shell of dust is released from rest, a singularity will be
produced at \(r=0\) when all the matter simultaneously reaches the origin. However,
if one perturbs the shell away from spherical symmetry or gives it some rotation, then
no such singularity will occur. Thus, one might suspect that in general relativity,
nonspherical collapse will not, in general, lead to a singularity. Similarly, if one
evolves a universe which is not exactly homogeneous and isotropic backward into
the past, perhaps instead of a ``big bang'' singularity, general relativity might predict
merely a nonsingular, high density phase which might then reexpand into the past.
If this occurred, then in general relativity one would have a viable nonsingular closed
universe model with infinitely many cycles of expansion, contraction, and ``bounce.''
Thus, it is of considerable interest to determine if the above predictions of
singularities are merely artifacts of the high degree of symmetry of the known exact
solutions, or whether singularities are a true, generic feature of solutions describing
cosmology and gravitational collapse in general relativity.

The purpose of this chapter is to discuss the singularity theorems of general
relativity. These theorems prove that singularities \emph{are} true, generic features of
cosmological and collapse solutions. Although they give very little information
concerning the detailed nature of these singularities, they show that models such as
the nonsingular ``bouncing universe'' are incompatible with general relativity
provided only that certain energy conditions are satisfied by matter and several other
conditions hold in the spacetime.

The prediction of singularities undoubtedly represents a breakdown of general
relativity in that its classical description of gravitation and matter cannot be expected
to remain valid at the extreme conditions expected near a spacetime singularity.
Indeed, on dimensional grounds one would expect a classical description of spacetime
structure to break down at scales characterized by the Planck length,

\[
  l_{\mathrm P}=(\hbar G/c^3)^{1/2}\simeq10^{-33}\,\mathrm{cm}
\]

(see chapter~14). Hence, classical general relativity certainly should break down at
or before the stage where it predicts spacetime curvatures of order

\[
  l_{\mathrm P}^{-2}.
\]

Although the singularity theorems do not prove that the singularities of classical
general relativity must involve unboundedly large curvature, they strongly suggest
the occurrence in cosmology and gravitational collapse of conditions in which
quantum or other effects which invalidate classical general relativity will play a
dominant role.

We begin our discussion of the singularity theorems in section~9.1 by attempting
to define the term ``singularity.'' Although it may be easy to recognize the ``big bang''
of the Robertson--Walker solutions or ``\(r=0\)'' of the Schwarzschild solution as
singularities, it is extremely difficult to give a satisfactory general notion of that
term. We provide motivation for the notion of timelike and null geodesic
incompleteness as a criterion for the presence of a singularity, although even this notion
has some unwanted features. It is this criterion which is used in the singularity
theorems.

The basic idea of the singularity theorems is as follows. As already shown in
section~3.3, geodesics are curves of extremal length. After deriving some useful
equations describing geodesic congruences in section~9.2, we obtain criteria in
section~9.3 for when a timelike geodesic fails to be a local maximum of ``length''
(i.e., proper time) between two points or between a point and a three-dimensional
surface. We obtain similar criteria for when a null geodesic fails to remain on the
boundary of the future of a point or two-dimensional surface. Using an inequality on
the Ricci tensor which follows from local positivity properties of the stress-energy
tensor---this being the only place in the analysis where Einstein's equation is used---
we show that in appropriate circumstances ``sufficiently long'' timelike geodesics
cannot be maximal length curves, and ``sufficiently long'' null geodesics cannot
remain on past or future boundaries. However, in section~9.4 we give global
arguments involving compactness of the spaces of causal curves (see section~8.3) to
establish existence of timelike curves of maximal length in globally hyperbolic
spacetimes. An analogous result (theorem~8.3.11) for null geodesics was already
given in chapter~8. The contradiction between these results (even if the spacetime is
not assumed to be globally hyperbolic) produces the singularity theorems of section
9.5: In appropriate circumstances ``sufficiently long'' timelike or null geodesics
cannot exist; i.e., one must have timelike or null geodesic incompleteness.
